\section{Python Installation}

\subsection{Installing Python 3.9+}

\textbf{Windows:}
\begin{lstlisting}[style=shell]
# Download from python.org and run installer
# Or use winget
winget install Python.Python.3.9
\end{lstlisting}

\textbf{macOS:}
\begin{lstlisting}[style=shell]
# Using Homebrew
brew install python@3.9
\end{lstlisting}

\textbf{Linux:}
\begin{lstlisting}[style=shell]
# Ubuntu/Debian
sudo apt update
sudo apt install python3.9 python3.9-venv python3-pip

# CentOS/RHEL
sudo yum install python39
\end{lstlisting}

\section{Docker Installation}

\subsection{Docker Desktop}

\textbf{Windows and macOS:}
\begin{enumerate}
    \item Download Docker Desktop from \url{https://www.docker.com/products/docker-desktop}
    \item Run the installer
    \item Start Docker Desktop
    \item Verify installation:
\end{enumerate}

\begin{lstlisting}[style=shell]
docker --version
docker run hello-world
\end{lstlisting}

\textbf{Linux:}
\begin{lstlisting}[style=shell]
# Install Docker Engine
curl -fsSL https://get.docker.com -o get-docker.sh
sudo sh get-docker.sh

# Add user to docker group
sudo usermod -aG docker $USER

# Start Docker
sudo systemctl start docker
sudo systemctl enable docker
\end{lstlisting}

\section{Kubernetes Setup}

\subsection{Minikube for Local Development}

\begin{lstlisting}[style=shell]
# Install minikube
# macOS
brew install minikube

# Linux
curl -LO https://storage.googleapis.com/minikube/releases/latest/minikube-linux-amd64
sudo install minikube-linux-amd64 /usr/local/bin/minikube

# Start minikube
minikube start --cpus 4 --memory 8192
\end{lstlisting}

\section{Cloud Provider Setup}

\subsection{AWS CLI}

\begin{lstlisting}[style=shell]
# Install AWS CLI
pip install awscli

# Configure credentials
aws configure
\end{lstlisting}

\subsection{Google Cloud SDK}

\begin{lstlisting}[style=shell]
# Install gcloud CLI
# Follow instructions at cloud.google.com/sdk/docs/install

# Initialize
gcloud init
\end{lstlisting}

\subsection{Azure CLI}

\begin{lstlisting}[style=shell]
# Install Azure CLI
pip install azure-cli

# Login
az login
\end{lstlisting}

\section{ML Frameworks Installation}

\subsection{PyTorch}

\begin{lstlisting}[style=shell]
# CPU version
pip install torch torchvision torchaudio

# GPU version (CUDA 11.8)
pip install torch torchvision torchaudio --index-url https://download.pytorch.org/whl/cu118
\end{lstlisting}

\subsection{TensorFlow}

\begin{lstlisting}[style=shell]
# CPU version
pip install tensorflow

# GPU version
pip install tensorflow[and-cuda]
\end{lstlisting}

\subsection{Scikit-learn and Data Science Stack}

\begin{lstlisting}[style=shell]
pip install numpy pandas scikit-learn matplotlib seaborn jupyter
\end{lstlisting}

\section{MLOps Tools Installation}

\subsection{MLflow}

\begin{lstlisting}[style=shell]
pip install mlflow

# Start MLflow server
mlflow server --backend-store-uri sqlite:///mlflow.db --default-artifact-root ./mlruns --host 0.0.0.0
\end{lstlisting}

\subsection{DVC}

\begin{lstlisting}[style=shell]
pip install dvc

# Initialize DVC
dvc init

# Configure remote storage
dvc remote add -d myremote s3://mybucket/dvc-storage
\end{lstlisting}

\subsection{Apache Airflow}

\begin{lstlisting}[style=shell]
pip install apache-airflow

# Initialize database
airflow db init

# Create admin user
airflow users create --username admin --firstname Admin --lastname User --role Admin --email admin@example.com

# Start webserver and scheduler
airflow webserver -p 8080
airflow scheduler
\end{lstlisting}

\section{GPU Drivers}

\subsection{NVIDIA Drivers and CUDA}

\textbf{Ubuntu/Debian:}
\begin{lstlisting}[style=shell]
# Install NVIDIA drivers
sudo apt install nvidia-driver-525

# Install CUDA Toolkit
wget https://developer.download.nvidia.com/compute/cuda/repos/ubuntu2204/x86_64/cuda-keyring_1.0-1_all.deb
sudo dpkg -i cuda-keyring_1.0-1_all.deb
sudo apt update
sudo apt install cuda

# Verify installation
nvidia-smi
nvcc --version
\end{lstlisting}

\section{Development Tools}

\subsection{VS Code Extensions for ML}

Recommended VS Code extensions:
\begin{itemize}
    \item Python
    \item Jupyter
    \item Docker
    \item Kubernetes
    \item YAML
    \item GitLens
    \item Remote - SSH
    \item Remote - Containers
\end{itemize}

\section{Complete Environment Setup Script}

\begin{lstlisting}[style=shell, caption=setup.sh]
#!/bin/bash

# Create project directory
mkdir -p ml_project
cd ml_project

# Create virtual environment
python -m venv venv
source venv/bin/activate

# Install dependencies
pip install --upgrade pip
pip install torch torchvision torchaudio
pip install tensorflow
pip install scikit-learn pandas numpy matplotlib seaborn
pip install mlflow dvc
pip install jupyter jupyterlab
pip install fastapi uvicorn
pip install pytest pytest-cov
pip install black flake8 mypy

# Initialize git
git init
curl https://raw.githubusercontent.com/github/gitignore/main/Python.gitignore -o .gitignore

# Initialize DVC
dvc init

# Create project structure
mkdir -p data/{raw,processed,features}
mkdir -p models
mkdir -p notebooks
mkdir -p src/{data,features,models,utils}
mkdir -p tests
mkdir -p config

# Create requirements.txt
pip freeze > requirements.txt

echo "Setup complete!"
\end{lstlisting}
